% =============================================================
% Aufgaben-Umgebungen.
%
% Drei Komfort-Umgebungen für die drei Aufgabentypen:
%   \begin{bleistiftuebung}{N}{Titel}  ...  \end{bleistiftuebung}
%   \begin{pythoneinheit} {N}{Titel}  ...  \end{pythoneinheit}
%   \begin{werkzeugcheck}    {Titel}  ...  \end{werkzeugcheck}
%
% Jede Umgebung erzeugt:
%   - eine farbige tcolorbox um den Aufgabentext, mit einem 2,5-pt-Strich
%     am AUSSEN-Rand der Seite (recto = rechts, verso = links),
%   - eine Marginalie auf derselben Außenseite mit Symbol und
%     Kapitälchen-Bezeichnung in der Aufgaben-Farbe.
%
% Die Bausteine stammen aus doku/LAYOUT_BRIEFING.md, dort steht auch
% die Begründung für jede Designentscheidung.
% =============================================================

% --- aufgabebox: tcolorbox mit page-aware Strich am Außenrand ---
% Eine einzige Box, kein Recto/Verso-Split mehr. Der overlay-Hook wird
% bei jedem Page-Part neu ausgewertet (auch wenn die Aufgabe per
% breakable über Seiten geht), prüft die aktuelle Seitenparität und
% zeichnet den Strich auf der jeweiligen Außenseite.
%
% left/right=8pt sorgt für Innenabstand zum Strich auf beiden Seiten;
% der Text bleibt immer mittig in der Box.
\newtcolorbox{aufgabebox}[1]{%
  enhanced, breakable,
  colback=white, colframe=white,
  boxrule=0pt, arc=0pt, outer arc=0pt,
  boxsep=0pt,
  left=8pt, right=8pt, top=4pt, bottom=4pt,
  before skip=0.6em, after skip=0.6em,
  overlay={%
    \checkoddpage
    \ifoddpage
      \draw[line width=2.5pt, color=#1]
        (frame.north east) -- (frame.south east);%
    \else
      \draw[line width=2.5pt, color=#1]
        (frame.north west) -- (frame.south west);%
    \fi
  },
}

% --- Marginalien-Inhalt für Recto und Verso ---
% Auf Recto-Seiten (Marginalspalte rechts): linksbündig zum Strich.
% Auf Verso-Seiten (Marginalspalte links):  rechtsbündig zum Strich.
% Das \vspace*{12pt} kompensiert das `before skip` und das `top` der
% tcolorbox, damit das Symbol mit der ersten Zeile der Aufgabenbox fluchtet.
\newcommand{\markercontentRecto}[4]{%
  \vspace*{12pt}\raggedright\color{#4}%
  {\markerfont\fontsize{24}{26}\selectfont #1}\\[0.05em]
  {\small\textsc{#2}}\\
  {\small\textsc{#3}}%
}
\newcommand{\markercontentVerso}[4]{%
  \vspace*{12pt}\raggedleft\color{#4}%
  {\markerfont\fontsize{24}{26}\selectfont #1}\\[0.05em]
  {\small\textsc{#2}}\\
  {\small\textsc{#3}}%
}

% --- Aufgaben-Umgebung ---
% Marginpar wird mit beiden Inhaltsvarianten aufgerufen (recto/verso);
% LaTeX wählt selbst die passende, je nachdem auf welcher Seite die
% Aufgabe beginnt. Der page-aware Strich entsteht im Overlay der
% aufgabebox und wechselt automatisch beim Seitenumbruch mit.
\newenvironment{aufgabe}[4]%
  {% #1 = Symbol, #2 = Label-Zeile-1, #3 = Label-Zeile-2, #4 = Farbe
   \marginpar[\markercontentVerso{#1}{#2}{#3}{#4}]%
              {\markercontentRecto{#1}{#2}{#3}{#4}}%
   \begin{aufgabebox}{#4}\ignorespaces}
  {\end{aufgabebox}}

% --- Komfort-Wrapper für die drei Aufgabentypen ---
% Konvention: der erste Teil der Aufgabe ist der Titel (fett), gefolgt
% vom Aufgabentext. Der Titel wird hier als Argument übergeben.
%
% Wir nutzen {\bfseries ...\par} statt \textbf{...\par}, weil \textbf
% als Inline-Befehl keinen Absatzumbruch (\par) im Argument verträgt
% ("Paragraph ended before \text@command was complete").
\newenvironment{bleistiftuebung}[2]%
  {\begin{aufgabe}{✏}{Bleistift-}{Übung #1}{cBleistift}{\bfseries #2.\par}}
  {\end{aufgabe}}

\newenvironment{pythoneinheit}[2]%
  {\begin{aufgabe}{⌨}{Python-}{Einheit #1}{cPython}{\bfseries #2.\par}}
  {\end{aufgabe}}

\newenvironment{werkzeugcheck}[1]%
  {\begin{aufgabe}{⚙}{Werkzeug-}{Check}{cWerkzeug}{\bfseries #1.\par}}
  {\end{aufgabe}}

% --- Code-Snippet-Einbindung ---
% \pythoncode{tagN/file.py} bindet eine Python-Datei vollständig ein.
% Pfad ist relativ zu latex/code/. Die .py-Dateien sind eigenständig
% lauffähig, sodass Greta sie 1:1 in Thonny verwenden kann.
\newcommand{\pythoncode}[1]{\inputminted{python}{code/#1}}

% \pythoncodeteil{tagN/file.py}{regionname} zeigt nur eine Region
% aus einer Datei. Region-Marker in der Quelle:
%     # region <name>
%     ... Code ...
%     # endregion
% Greta tippt die Datei sukzessive auf, im Buch zeigen wir pro
% Python-Einheit jeweils nur das aktuell relevante Stück.
%
% Implementierung: extract_region.py schneidet die Region zur
% Build-Zeit in eine temporäre Datei, die minted dann einbindet.
% Der Counter sorgt dafür, dass jede Einbindung ihre eigene Datei
% bekommt — nötig, weil minted das Snippet zur Build-Zeit cached.
\newcounter{snippetcntr}
\makeatletter
\newcommand{\pythoncodeteil}[2]{%
  \stepcounter{snippetcntr}%
  \protected@edef\@snippetfile{.snippets/snippet-\arabic{snippetcntr}.py}%
  \immediate\write18{python3 scripts/extract_region.py code/#1 #2 > \@snippetfile}%
  \inputminted{python}{\@snippetfile}%
}
\makeatother
